% =====================================================================
%  RESUMEN EJECUTIVO (castellano / valenciano / inglés)
%  NOTA: la versión en valenciano debería ser revisada por un hablante nativo.
% =====================================================================
\chapter*{Resumen ejecutivo}
\addcontentsline{toc}{chapter}{Resumen ejecutivo}
\markboth{Resumen ejecutivo}{Resumen ejecutivo}

\section*{Resumen}

Este Trabajo Fin de Grado analiza, mediante técnicas de procesamiento del lenguaje natural
(PLN), cómo las grandes empresas europeas del índice STOXX Europe 600 comunican su estrategia
de sostenibilidad bajo el nuevo marco normativo de la CSRD y los estándares ESRS, con el
objetivo de identificar patrones comunicativos y señales textuales de \textit{greenwashing}.
A partir de un corpus propio de 586 informes corporativos de 196 empresas, correspondientes a
los ejercicios 2022, 2023 y 2024, se aíslan las subsecciones de sostenibilidad y se les aplica
una batería de técnicas: diccionarios de tono y de cobertura ESRS, dos métodos de modelado de
temas (LDA y BERTopic) y modelos de lenguaje especializados (FinBERT y ClimateBERT). Con sus
resultados se construye un índice de \textit{greenwashing} textual y se realizan diversos
contrastes estadísticos. En coherencia con su objeto, el análisis prescinde por completo de
calificaciones ESG de terceros y se basa únicamente en el texto de los informes. Los resultados
muestran que el discurso se organiza en torno al clima y a las personas; que difiere de forma
muy significativa por sector y país —el sector financiero y el tecnológico encabezan el lenguaje
vago—; que la hipótesis simple de que el optimismo encubre vaguedad no se sostiene; y, sobre
todo, que el tránsito de la NFRD a la CSRD transforma profundamente el reporting: los informes
casi triplican su extensión, adoptan un tono menos optimista, incorporan más lenguaje de riesgo
y elevan su índice de \textit{greenwashing}, con la doble materialidad como emblema del nuevo
régimen. El trabajo aporta, además, un cuadro de mando interactivo que permite explorar estos
hallazgos.

\medskip
\noindent\textbf{Palabras clave:} \textit{greenwashing}; CSRD; ESRS; procesamiento del lenguaje
natural; informes de sostenibilidad; STOXX Europe 600; análisis de contenido; ClimateBERT.

\section*{Resum}

Aquest Treball Fi de Grau analitza, mitjançant tècniques de processament del llenguatge natural
(PLN), com les grans empreses europees de l'índex STOXX Europe 600 comuniquen la seua estratègia
de sostenibilitat sota el nou marc normatiu de la CSRD i els estàndards ESRS, amb l'objectiu
d'identificar patrons comunicatius i senyals textuals de \textit{greenwashing}. A partir d'un
corpus propi de 586 informes corporatius de 196 empreses, corresponents als exercicis 2022, 2023
i 2024, s'aïllen les subseccions de sostenibilitat i se'ls apliquen diverses tècniques:
diccionaris de to i de cobertura ESRS, dos mètodes de modelatge de temes (LDA i BERTopic) i
models de llenguatge especialitzats (FinBERT i ClimateBERT). Amb els seus resultats es construeix
un índex de \textit{greenwashing} textual i es realitzen diversos contrastos estadístics. En
coherència amb el seu objecte, l'anàlisi prescindeix completament de les qualificacions ESG de
tercers i es basa únicament en el text dels informes. Els resultats mostren que el discurs
s'organitza al voltant del clima i de les persones; que difereix de manera molt significativa per
sector i país —el sector financer i el tecnològic encapçalen el llenguatge vague—; que la
hipòtesi simple que l'optimisme amaga vaguetat no se sosté; i, sobretot, que el pas de la NFRD a
la CSRD transforma profundament la informació: els informes quasi tripliquen la seua extensió,
adopten un to menys optimista, incorporen més llenguatge de risc i eleven el seu índex de
\textit{greenwashing}, amb la doble materialitat com a emblema del nou règim. El treball aporta,
a més, un quadre de comandament interactiu que permet explorar aquestes troballes.

\medskip
\noindent\textbf{Paraules clau:} \textit{greenwashing}; CSRD; ESRS; processament del llenguatge
natural; informes de sostenibilitat; STOXX Europe 600; anàlisi de contingut; ClimateBERT.

\section*{Abstract}

This Bachelor's Thesis analyses, by means of natural language processing (NLP) techniques, how
large European companies in the STOXX Europe 600 index communicate their sustainability strategy
under the new regulatory framework of the CSRD and the ESRS standards, with the aim of
identifying communicative patterns and textual signals of greenwashing. Drawing on a purpose-built
corpus of 586 corporate reports from 196 companies, covering the 2022, 2023 and 2024 financial
years, the sustainability subsections are isolated and analysed with a range of techniques: tone
and ESRS-coverage dictionaries, two topic-modelling methods (LDA and BERTopic) and specialised
language models (FinBERT and ClimateBERT). Their outputs feed a textual greenwashing index and a
set of statistical tests. Consistent with its object of study, the analysis deliberately avoids
any third-party ESG ratings and relies solely on the text of the reports. The findings show that
the discourse revolves around climate and people; that it differs very significantly across
sectors and countries —with the financial and technology sectors leading the use of vague
language—; that the simple hypothesis whereby optimism conceals vagueness does not hold; and,
above all, that the transition from the NFRD to the CSRD profoundly reshapes reporting: reports
almost triple in length, adopt a markedly less optimistic tone, incorporate more risk language
and raise their greenwashing index, with double materiality as the emblem of the new regime. The
work also delivers an interactive dashboard that makes it possible to explore these findings.

\medskip
\noindent\textbf{Keywords:} greenwashing; CSRD; ESRS; natural language processing; sustainability
reports; STOXX Europe 600; content analysis; ClimateBERT.

\clearpage
